\chapter*{Introduction}
\addcontentsline{toc}{chapter}{Introduction}
\markboth{INTRODUCTION}{INTRODUCTION}

\section*{Who this book is for}

This book is written for the working backend engineer who is fluent in a typed,
async-first ecosystem --- most likely .NET --- and who is now being asked to put
a Python agent into production.

It assumes you have shipped services. You know what a connection pool is for,
why idempotency matters on a retried request, what a cancellation token does,
and roughly what you would look at first when a service degrades under load. It
assumes you have written \code{async}/\code{await} somewhere, in some language.

It does \emph{not} assume you know Python well. Chapter~\ref{ch:python} exists
precisely because the language is the smallest part of the gap and the part most
likely to make an otherwise strong engineer write code a reviewer dislikes.
It does not assume you have built an agent before, and it does not assume you
have read a paper about large language models.

It assumes one more thing, which is less common: that you are sceptical. The
overwhelming majority of agent material published since 2023 demonstrates that a
thing is \emph{possible} without asking whether it is \emph{reliable},
\emph{observable}, or \emph{affordable}. A notebook that calls three tools and
prints a paragraph proves nothing about a service that must hold two hundred
concurrent conversations, survive a pod restart mid-conversation, wait three days
for a human to approve an irreversible action, and stay inside a budget. This
book is organised around the second kind of system.

\section*{What this book is not}

\begin{itemize}[leftmargin=1.4em]
  \item It is not an introduction to large language models.
        Chapter~\ref{ch:landscape} assumes you know roughly what a token, a
        context window, and a tool call are.
  \item It is not a Python tutorial. Chapter~\ref{ch:python} covers what a
        senior engineer from another ecosystem actually gets wrong, which is a
        much smaller and more specific set of things than a tutorial covers.
  \item It is not prompt engineering. There are exactly two places where the
        wording of a prompt is treated as an engineering concern: tool
        descriptions, because they measurably change tool-selection accuracy,
        and judge prompts, because a miscalibrated judge silently corrupts every
        number downstream of it.
  \item It is not a LangChain advocacy document. Chapter~\ref{ch:ecosystem} sets
        out the alternatives, the cases where no framework is the right answer,
        and the criticisms of LangChain that are actually well founded --- as
        distinct from the ones that were true in 2023 and have been repeated ever
        since.
\end{itemize}

\section*{Why this book exists}

On 22~October~2025, LangChain and LangGraph both reached 1.0, and the
relationship between them changed. \pkg{langchain} stopped being a large library
of chains and became a thin, opinionated layer over \pkg{langgraph}'s execution
runtime. \api{AgentExecutor}, \api{initialize\_agent} and \api{LLMChain} --- the
three things every tutorial written before that date is built on --- were
retired. Retrievers, indexing and the chain abstractions moved to
\pkg{langchain-classic}, which is in maintenance.

The practical consequence is unusually severe: \emph{most of the LangChain
material on the internet describes an architecture that no longer exists}. Search
results, blog posts, conference talks and a very large amount of training data
all predate the split. It is entirely possible to follow a well-written,
well-rated tutorial from 2024 and end up with code that no current maintainer
would accept.

That is the gap this book is written into. The 1.0 line carries a commitment to
no breaking changes before 2.0, so for the first time the material is worth
writing down carefully rather than provisionally.

\section*{How the book is organised}

Five parts, seventeen chapters.

\textbf{Part I --- Foundations}
(Chapters~\ref{ch:landscape}--\ref{ch:cancellation}) covers the ground the rest
of the book stands on: the state of the ecosystem and what died in 1.0, the
Python a senior engineer needs and the Python idioms they tend to get wrong, and
then two chapters on \code{asyncio}. Two, not one, because the job advertisement
that sent you here said ``async Python'' for a reason: an agent is a workload in
which better than ninety per cent of wall-clock time is spent waiting on a
network, and every interesting failure mode in
Chapter~\ref{ch:cancellation} --- a swallowed cancellation, an unbounded queue, a
blocking call in a coroutine --- shows up in production rather than in a notebook.

\textbf{Part II --- LangChain 1.x} (Chapters~\ref{ch:models}--\ref{ch:agents})
covers the layer most people mean when they say ``LangChain'': models, messages
and the \api{Runnable} interface; tools and structured output; and then
\api{create\_agent} and the middleware system that is 1.0's flagship addition.

\textbf{Part III --- LangGraph} (Chapters~\ref{ch:stategraph}--\ref{ch:multi-agent})
covers the runtime underneath, and is the longest and most important part of the
book. State and reducers, the Pregel execution model, persistence and durable
execution, interrupts and time travel, streaming, and multi-agent topologies with
an honest account of what they cost.

\textbf{Part IV --- Production} (Chapters~\ref{ch:context}--\ref{ch:security})
covers what you need before any of the above is allowed near a user: managing the
context window as the scarce resource it is, serving agents over HTTP with
streaming and concurrency control, observability and evaluation, and then
security, cost and deployment.

\textbf{Part V --- Perspective} (Chapters~\ref{ch:ecosystem}--\ref{ch:capstone})
steps outside the framework --- MCP, Deep Agents, the competing SDKs, the cases
for using nothing at all --- and then reviews the finished mini-project as a set
of decisions rather than a set of files.

Chapters are independent enough to read out of order after Part~I, with two
exceptions: Chapters~\ref{ch:stategraph} and~\ref{ch:persistence} must be read in
sequence, and Chapter~\ref{ch:cancellation} assumes
Chapter~\ref{ch:asyncio}.

\section*{The mini-project}

The book builds one system across its length, in stages. It is called
\textbf{Ops Copilot}: an internal operations assistant, exposed as a FastAPI
service, that answers questions about orders and incidents by querying real
systems, and that can take one irreversible action --- sending an email --- only
behind a human approval gate.

\begin{center}
\code{github.com/konradcinkusz/llm-book-mini-project}
\end{center}

The choice is deliberate. It is small enough to hold in your head and large
enough to be genuinely hard in the places that matter: it needs concurrency, it
needs persistence that survives a restart, it needs a human in the loop, it needs
to stream, it needs to not send the same email twice when a node replays, and it
needs to cost a defensible amount of money per conversation.

Each chapter that adds to it opens a \emph{Mini-project} callout naming the stage
it builds and what that stage must demonstrate. Every stage has tests, and the
tests run without a single real model call.

\section*{Conventions}

Six kinds of callout appear throughout.

\begin{note}
Background, context, or a design rationale. Skippable on a first pass.
\end{note}

\begin{warning}
A failure mode that will cost you an afternoon if you meet it unprepared.
Not skippable.
\end{warning}

\begin{versionbox}
Applies to a surface that is new, experimental, or likely to move between minor
versions. The LangChain ecosystem releases weekly; check the release notes before
relying on anything marked this way.
\end{versionbox}

\begin{verifybox}
A listing transcribed from documentation or release notes rather than executed by
the author against the pinned versions. Run it before you copy it into anything
that matters.
\end{verifybox}

\begin{dotnetbox}
A translation of the surrounding idea into .NET terms. These carry no content of
their own --- skip every one of them and you lose nothing but a shortcut.
\end{dotnetbox}

\begin{projectbox}
The stage of the mini-project that this section builds, and what it must
demonstrate when it is finished.
\end{projectbox}

\noindent
Inline, \pkg{langgraph} denotes a distribution on PyPI, \api{create\_agent} an
API member, and \code{uv sync} a literal command or file path.

\section*{On versions, and why this book will rot}

LangChain is not one package on one version line. It is a family of distributions
that release independently and frequently: at the time of writing,
\pkg{langchain}~\lcver{}, \pkg{langchain-core}~\lccorever{},
\pkg{langgraph}~\lgver{}, and a separate version line for each provider
integration and each checkpointer backend. A version number quoted for one of
them tells you very little about the others.

Every listing in this book is written against the versions on the title page, and
the mini-project pins them exactly in its lockfile. A CI job compares those pins
against PyPI on every build, so a stale pin is visible rather than discovered by
a reader.

The honest position: the conceptual material --- the execution model, reducers,
the replay semantics of a node, the economics of a supervisor --- will outlive
many minor versions, because those are load-bearing and 1.0 committed to them.
The exact import paths and keyword arguments will not. Treat the listings as
something you run, not as scripture. Where the book has not run something, it
says so in a box rather than rounding up.

\section*{Environment assumptions}

\begin{itemize}[leftmargin=1.4em]
  \item \pyver{} (\pymin{} is the floor LangChain 1.x supports; the book uses
        \code{asyncio.TaskGroup} and \code{asyncio.timeout} throughout, which
        need 3.11+).
  \item \code{uv} for environment and dependency management.
        Chapter~\ref{ch:python} explains why, and what to do if your organisation
        has standardised on something else.
  \item Docker, for Postgres from Chapter~\ref{ch:persistence} onwards.
  \item An API key for at least one model provider with tool-calling support.
        The book is provider-neutral; the mini-project runs against any of them.
  \item Optional: a LangSmith account for Chapter~\ref{ch:observability}. That
        chapter also covers the self-hosted alternatives, because not every
        client will accept a SaaS tracing backend.
\end{itemize}

Chapter~\ref{ch:python} walks the setup end to end and assumes none of the above
is installed.

\vspace{1.5em}
\noindent\rule{\linewidth}{0.4pt}
\vspace{0.5em}

\noindent\emph{Konrad Cinkusz\\ \today}

\cleardoublepage
